\chapter{Discussion}
\label{ch:discussion}

% ============================================================================
% This chapter discusses results, limitations, and implications
% ============================================================================

\section{Overview}
\label{sec:discussion_overview}

This chapter discusses the experimental results, compares findings with prior work, analyzes limitations, and discusses implications for sleep stage classification research.

\section{Model Performance}
\label{sec:performance_discussion}

\subsection{XGBoost vs Random Forest}
\label{subsec:xgb_vs_rf}

% TODO: After results, discuss:
% - Which model performed better and why?
% - Are differences statistically significant?
% - What are trade-offs (accuracy vs training time, interpretability)?

Our results show that [XGBoost/Random Forest] achieved superior performance with [X\%] accuracy compared to [Y\%] for [other model]. This [is/is not] consistent with prior work in sleep stage classification \cite{TODO_ADD_REF}.

Possible explanations for these results include:

\begin{itemize}
    \item \textbf{Gradient boosting:} XGBoost's sequential error correction may be better suited for class imbalance in sleep data
    \item \textbf{Regularization:} XGBoost's built-in L1/L2 regularization may prevent overfitting on minority classes (N1)
    \item \textbf{Feature interactions:} [Discuss which model captures feature interactions better]
\end{itemize}

\subsection{Subject-Independent Generalization}
\label{subsec:loso_generalization}

The LOSO cross-validation results demonstrate [good/moderate/poor] subject-independent generalization. [Discuss variability across subjects - some subjects much easier/harder than others?]

% TODO: Add figure showing per-subject accuracy distribution (box plot or histogram)

Factors affecting inter-subject variability may include:

\begin{itemize}
    \item Age and gender differences
    \item Sleep quality and architecture
    \item Individual EEG signal characteristics
    \item Recording quality and artifacts
\end{itemize}

\subsection{Per-Class Performance}
\label{subsec:perclass_discussion}

% TODO: Discuss confusion patterns
% - Which stages are most confused with each other?
% - Is N1 (rare class) poorly classified as expected?
% - Are transitions (W→N1, N1→N2) harder than stable periods?

Common confusion patterns observed:

\begin{itemize}
    \item \textbf{N1 vs Wake:} N1 is often confused with Wake due to similar EEG characteristics during drowsiness
    \item \textbf{N1 vs N2:} Transition between light sleep stages is ambiguous
    \item \textbf{REM vs Wake:} Paradoxical sleep has EEG similar to wakefulness
\end{itemize}

These patterns are consistent with clinical sleep scoring challenges \cite{TODO_ADD_REF}.

\section{Feature Selection Analysis}
\label{sec:feature_discussion}

\subsection{Impact of Feature Count}
\label{subsec:feature_count}

% TODO: Discuss top-K = 30 vs 50 vs 75
% - Does performance plateau?
% - Optimal trade-off between dimensionality and information?

Our results [show/do not show] diminishing returns with increasing feature count. Top-K = [VALUE] provides the best balance between dimensionality and performance.

\subsection{Important Features}
\label{subsec:important_features_discussion}

% TODO: Analyze which features were most selected
% - Are frequency-domain features more important than time-domain? (Expected for sleep)
% - Which EEG channels are most informative? (Central channels? Frontal?)
% - Are delta and theta bands dominant? (Expected for deep sleep and REM)

The most frequently selected features include [list top features]. This suggests that [discuss what this tells us about sleep stage signatures].

Frequency-domain features [dominate/complement] time-domain features, consistent with the understanding that sleep stages are characterized by distinct spectral patterns \cite{TODO_ADD_REF}.

\section{Computational Efficiency}
\label{sec:efficiency_discussion}

\subsection{Caching System Impact}
\label{subsec:caching_impact}

The two-tier caching system achieved [X\%] feature cache hit rate and [Y\%] LOSO model cache hit rate, reducing total experiment time by [Z hours/days].

% TODO: Compare with naive implementation
% - Without caching: X hours for full experiment
% - With caching (after initial run): Y hours for parameter sweep
% - Speedup factor: X/Y

This demonstrates the practical value of fingerprint-based caching for iterative machine learning experimentation.

\subsection{Parallelization}
\label{subsec:parallelization}

Epoch-level parallelization in feature extraction provided [X×] speedup on [N-core] CPU. Subject-level parallelization was not implemented to avoid memory overhead and cache contention, as caching already makes subject processing nearly instant on cache hits.

\section{Comparison with Prior Work}
\label{sec:comparison}

% TODO: Compare with literature
% - How does our accuracy/kappa compare to other BOAS studies?
% - How does our accuracy compare to other sleep staging methods?
% - Is LOSO evaluation more stringent than random K-fold?

Table~\ref{tab:literature_comparison} compares our results with prior work.

\begin{table}[htbp]
    \centering
    \caption{Comparison with Prior Sleep Stage Classification Studies}
    \label{tab:literature_comparison}
    \begin{tabular}{llccc}
        \toprule
        \textbf{Study} & \textbf{Method} & \textbf{Dataset} & \textbf{CV} & \textbf{Accuracy} \\
        \midrule
        % TODO: Fill in literature comparison
        Author 2023 \cite{TODO} & CNN & Sleep-EDF & 10-fold & XX.X\% \\
        Author 2022 \cite{TODO} & LSTM & MASS & Random & XX.X\% \\
        Author 2021 \cite{TODO} & Random Forest & Custom & LOSO & XX.X\% \\
        \midrule
        \textbf{This work} & XGBoost & BOAS & LOSO & XX.X\% \\
        \bottomrule
    \end{tabular}
\end{table}

Our results [are comparable to / exceed / fall short of] state-of-the-art methods, considering that:

\begin{itemize}
    \item LOSO is more stringent than random K-fold cross-validation
    \item Different datasets have different difficulty levels
    \item Feature-based methods are more interpretable than deep learning
\end{itemize}

\section{Limitations}
\label{sec:limitations}

This work has several limitations:

\subsection{Dataset Limitations}
\label{subsec:dataset_limitations}

\begin{itemize}
    \item \textbf{Single dataset:} Results are specific to BOAS; generalization to other datasets is uncertain
    \item \textbf{Healthy subjects:} BOAS contains only healthy sleepers; clinical populations (sleep disorders) may have different characteristics
    \item \textbf{Age range:} [Check BOAS demographics - limited age range?]
\end{itemize}

\subsection{Methodological Limitations}
\label{subsec:method_limitations}

\begin{itemize}
    \item \textbf{Handcrafted features:} Feature engineering requires domain expertise; deep learning may discover better representations
    \item \textbf{Hyperparameter tuning:} Limited hyperparameter search due to computational constraints
    \item \textbf{Class imbalance:} No special handling (e.g., SMOTE, class weights) for rare stages like N1
    \item \textbf{Single-epoch classification:} No temporal context from neighboring epochs (sequence models could help)
\end{itemize}

\subsection{Implementation Limitations}
\label{subsec:impl_limitations}

\begin{itemize}
    \item \textbf{Sequential LOSO:} Folds are processed sequentially; could parallelize across folds (but cache makes this less critical)
    \item \textbf{Memory constraints:} Feature extraction batch size limited by RAM
    \item \textbf{Storage overhead:} Cache requires [X GB] storage for 128 subjects
\end{itemize}

\section{Implications}
\label{sec:implications}

\subsection{For Sleep Research}
\label{subsec:research_implications}

This work demonstrates that:

\begin{itemize}
    \item Feature-based machine learning remains competitive for sleep stage classification
    \item LOSO cross-validation is essential for honest subject-independent evaluation
    \item Computational efficiency (via caching) enables extensive experimentation
\end{itemize}

\subsection{For Clinical Applications}
\label{subsec:clinical_implications}

The pipeline could be adapted for:

\begin{itemize}
    \item Automated sleep stage scoring in clinical sleep labs
    \item Home sleep monitoring with portable EEG devices
    \item Large-scale epidemiological sleep studies
\end{itemize}

However, clinical deployment requires:
\begin{itemize}
    \item Validation on clinical populations (sleep disorders)
    \item Regulatory approval (medical device certification)
    \item Integration with clinical workflows
\end{itemize}

\subsection{For Machine Learning}
\label{subsec:ml_implications}

The two-tier caching system with fingerprint-based invalidation is a general pattern applicable to other machine learning domains with expensive preprocessing and iterative experimentation.

\section{Summary}
\label{sec:discussion_summary}

This chapter discussed model performance, feature selection, computational efficiency, comparison with prior work, limitations, and implications. [Summarize key takeaways after results are available.]
